% =====================================================================
%  Chapitre 4 — Sprint 2 : Éditeur de modèles multi-canal
% =====================================================================
\chapter{Sprint~2 : Éditeur de modèles multi-canal}
\label{chap:sprint2}

\section*{Introduction}
Ce deuxième sprint constitue le c\oe{}ur métier de la plateforme. Il met en
\oe{}uvre l'éditeur de modèles et, en particulier, l'éditeur d'e-mail par
glisser-déposer qui en est le produit phare, ainsi que la prise en charge
des cinq canaux (e-mail, facture, contrat, SMS, RCS) et la bibliothèque de
modèles. Après avoir énoncé les objectifs du sprint, nous détaillons les
cas d'utilisation, les diagrammes de séquence et la réalisation.

\section{Objectifs du sprint}
Ce sprint couvre les epics~4 à~6 du backlog. Le
tableau~\ref{tab:backlog-s2} en présente les user stories.

\begin{table}[H]
\centering
\caption{Backlog du sprint~2}
\label{tab:backlog-s2}
\small
\begin{tabularx}{\linewidth}{clX}
\toprule
\textbf{ID} & \textbf{User story} & \textbf{Description} \\
\midrule
4.1 & Éditer un e-mail & En tant qu'éditeur, je veux composer un e-mail en glissant-déposant des blocs sur un canevas. \\
4.2 & Prévisualiser et coder & En tant qu'éditeur, je veux prévisualiser le rendu en direct et accéder à l'éditeur de code. \\
4.3 & Exporter & En tant qu'éditeur, je veux exporter mon modèle en HTML responsive et en PDF. \\
5.1 & Facture & En tant qu'éditeur, je veux construire une facture structurée avec variables. \\
5.2 & Contrat & En tant qu'éditeur, je veux rédiger un contrat en texte riche paginé. \\
5.3 & SMS & En tant qu'éditeur, je veux composer un SMS avec des variables de personnalisation. \\
5.4 & RCS & En tant qu'éditeur, je veux construire un message RCS (carte, carrousel, boutons) avec aperçu téléphone. \\
6.1 & Favoris et duplication & En tant qu'éditeur, je veux mettre des modèles en favoris et les dupliquer. \\
6.2 & Galerie prédéfinie & En tant que membre marketing, je veux constituer une galerie de modèles prédéfinis par catégorie. \\
\bottomrule
\end{tabularx}
\end{table}

\section{Raffinement et description des cas d'utilisation}

\subsection{Cas d'utilisation « Gérer les modèles »}
La figure~\ref{fig:uc-templates} regroupe les cas d'utilisation liés à la
gestion des modèles : création, édition, prévisualisation, export,
duplication, mise en favori et gestion de la galerie prédéfinie.

% [FIGURE : diagramme de cas d'utilisation « Gérer les modèles »]
\figplaceholder{Cas d'utilisation « Gérer les modèles »}{fig:uc-templates}{7.5cm}

Le tableau~\ref{tab:uc-create-template} décrit le scénario « Créer un
modèle d'e-mail ».

\begin{table}[H]
\centering
\caption{Description textuelle du cas d'utilisation « Créer un modèle d'e-mail »}
\label{tab:uc-create-template}
\small
\begin{tabularx}{\linewidth}{lX}
\toprule
\textbf{Cas d'utilisation} & Créer un modèle d'e-mail \\
\textbf{Acteur} & Éditeur \\
\textbf{Préconditions} & L'utilisateur est connecté et son plan autorise la création. \\
\textbf{Postconditions} & Un modèle d'e-mail est créé et enregistré pour son entreprise. \\
\midrule
\textbf{Scénario nominal} &
1.~L'éditeur choisit le canal « e-mail » et ouvre l'éditeur.\newline
2.~Il glisse-dépose des blocs sur le canevas et ajuste leur contenu et leur style.\newline
3.~Il prévisualise le rendu en direct.\newline
4.~Il enregistre le modèle.\newline
5.~Le système valide et persiste le modèle rattaché à son entreprise. \\
\midrule
\textbf{Scénario alternatif} &
5.a.~Si le quota du plan est atteint, le système bloque la création et propose une montée en gamme. \\
\bottomrule
\end{tabularx}
\end{table}

Le tableau~\ref{tab:uc-export} décrit le scénario « Exporter un modèle ».

\begin{table}[H]
\centering
\caption{Description textuelle du cas d'utilisation « Exporter un modèle »}
\label{tab:uc-export}
\small
\begin{tabularx}{\linewidth}{lX}
\toprule
\textbf{Cas d'utilisation} & Exporter un modèle \\
\textbf{Acteur} & Éditeur \\
\textbf{Préconditions} & Un modèle existe et est ouvert. \\
\textbf{Postconditions} & Un fichier HTML ou PDF est généré et téléchargé. \\
\midrule
\textbf{Scénario nominal} &
1.~L'éditeur demande l'export dans le format souhaité.\newline
2.~Le système compile le contenu (MJML vers HTML pour l'e-mail) et substitue les variables.\newline
3.~Le système produit le fichier et le met à disposition de l'utilisateur. \\
\midrule
\textbf{Scénario alternatif} &
2.a.~En cas d'erreur de compilation, le système en informe l'utilisateur. \\
\bottomrule
\end{tabularx}
\end{table}

\section{Diagrammes de séquence}

\subsection{Séquence « Créer un modèle »}
La figure~\ref{fig:seq-create-template} décrit les échanges entre l'éditeur,
la passerelle API et le service de modèles lors de la création d'un modèle,
depuis la saisie jusqu'à la persistance en base.

% [FIGURE : diagramme de séquence « Créer un modèle »]
\figplaceholder{Diagramme de séquence « Créer un modèle »}{fig:seq-create-template}{7cm}

\subsection{Séquence « Exporter un modèle »}
La figure~\ref{fig:seq-export} illustre le processus d'export : compilation
du contenu, substitution des variables et génération du fichier.

% [FIGURE : diagramme de séquence « Exporter un modèle »]
\figplaceholder{Diagramme de séquence « Exporter un modèle »}{fig:seq-export}{7cm}

\section{Réalisation}

\subsection{Service de modèles et modèle de données par blocs}
Le service de modèles, développé avec NestJS selon le patron CQRS, expose en
gRPC l'ensemble des opérations du cycle de vie d'un modèle : création,
modification, suppression, consultation, listage, rendu, duplication,
gestion des favoris et récupération des modèles populaires. Les données sont
persistées dans la base \texttt{templates\_db}~; les médias (images) sont
stockés dans MinIO.

Chaque modèle est caractérisé par un \textbf{type de canal} — e-mail (1),
facture (2), contrat (3), SMS (4) ou RCS (5) — et par un contenu dont la
forme dépend du canal. Pour l'e-mail, le contenu suit un
\textbf{modèle de blocs} structuré : un modèle est composé de rangées
(\emph{rows}), chaque rangée pouvant comporter jusqu'à quatre colonnes, et
chaque colonne contenant des blocs. Onze types de blocs sont disponibles :
titre, texte, image, vidéo, bouton, séparateur, tableau, signature, réseaux
sociaux, menu et liste à icônes. Ce modèle de données présente une propriété
essentielle : il est \textbf{identique} à celui que produit l'assistant
d'IA (chapitre~\ref{chap:sprint3}), de sorte que l'édition manuelle et
l'édition assistée deviennent parfaitement interchangeables.

\subsection{Liste des modèles et bibliothèque}
La figure~\ref{fig:ui-templates-list} présente la liste des modèles de
l'entreprise, avec aperçu, mise en favori, duplication et suppression. Une
page dédiée regroupe les favoris, et une galerie expose les modèles
prédéfinis maintenus par les membres marketing, organisés par catégorie.

% [FIGURE : capture — liste des modèles]
\figplaceholder{Liste des modèles}{fig:ui-templates-list}{6cm}

\subsection{Éditeur d'e-mail visuel}
L'éditeur d'e-mail (figure~\ref{fig:ui-email-editor}) est le composant phare
de la plateforme. Il s'organise autour d'un canevas central sur lequel
l'utilisateur glisse-dépose des blocs, d'un panneau de blocs à gauche et
d'un panneau de style à droite permettant d'ajuster les propriétés (couleur,
espacement, alignement, largeurs de colonnes, styles de section). Le
glisser-déposer s'appuie sur la bibliothèque dnd-kit. Un aperçu de l'e-mail
est rendu en direct, et un éditeur de code (Monaco) est mis à disposition
des utilisateurs avancés. À l'export, le modèle est compilé en HTML
responsive au moyen de MJML, ou en PDF.

% [FIGURE : capture — éditeur d'e-mail (canevas, panneaux blocs et style)]
\figplaceholder{Éditeur d'e-mail visuel}{fig:ui-email-editor}{7cm}

\subsection{Éditeurs de facture, de contrat et de SMS}
Les canaux transactionnels et texte disposent chacun d'un éditeur adapté.
L'éditeur de \textbf{facture} propose un constructeur structuré avec
variables de personnalisation et génération PDF. L'éditeur de
\textbf{contrat} repose sur un composant de texte riche (Tiptap) avec
pagination et export PDF. L'éditeur de \textbf{SMS} offre une saisie simple
intégrant des variables du type \texttt{\{\{prénom\}\}}. La
figure~\ref{fig:ui-facture-contrat} en présente un aperçu.

% [FIGURE : captures — éditeurs facture / contrat / SMS]
\figplaceholder{Éditeurs de facture, de contrat et de SMS}{fig:ui-facture-contrat}{6cm}

\subsection{Éditeur RCS et aperçu téléphone}
Le canal RCS (\emph{Rich Communication Services}) a été ajouté comme
cinquième type de modèle. Son éditeur (figure~\ref{fig:ui-rcs}) associe un
formulaire de composition à un \textbf{aperçu téléphone en direct}. Il
permet de construire un message texte, une carte enrichie, un carrousel de
cartes, ainsi que des puces de suggestion (répondre, ouvrir une URL,
appeler). Le message est stocké sous la forme d'une charge utile JSON,
alignée sur la représentation RBM, dans le champ de contenu du modèle — de
la même manière que le MJML sert de contenu à l'e-mail. La plateforme se
concentre ici sur la construction, le stockage et le rendu avec substitution
de variables~; l'envoi effectif via un opérateur RCS constitue une
perspective d'évolution.

% [FIGURE : capture — éditeur RCS et aperçu téléphone]
\figplaceholder{Éditeur RCS avec aperçu téléphone}{fig:ui-rcs}{6cm}

\section*{Conclusion}
Ce sprint a livré le c\oe{}ur fonctionnel de la plateforme : un éditeur
d'e-mail visuel complet reposant sur un modèle de blocs structuré, la prise
en charge des cinq canaux de communication, l'export HTML et PDF, ainsi que
la bibliothèque de modèles (favoris, duplication, galerie prédéfinie). Le
caractère structuré et partagé du modèle de blocs prépare directement le
sprint suivant, consacré à l'assistant d'intelligence artificielle, qui
produira et modifiera ces mêmes blocs.
